\chapter{Problem Definition and State of the Art}
\label{chap:stateoftheart}

\section{Introduction}
\label{sec:intro}

Scoring targets is an essential task in sport shooting, and a variety of systems have been developed
over the years to replace the time-consuming manual scoring. Common technologies include
triangulation of acoustic signals, laser-beam interruption and impact sensing. These technologies
represent high-accuracy but high-cost solutions, and are well-suited for professional events. Many
amateur sport-shooting clubs are interested in low-cost, real-time scoring systems that do not
require the same level of precision. Computer-vision-based systems offer a low-cost alternative
capable of detecting and scoring paper targets.

This section reviews the state of the art in computer-vision-based target scoring systems. This includes
classical image processing methods, and more recent deep learning approaches.

\section{Physical and Regulatory Requirements of ISSF Paper Targets}\label{sec:requirements}

ISSF defines the dimensions and scoring rules for paper targets used in 10\,m air rifle and 10\,m air pistol events.

\begin{figure}[ht]
\centering

\begin{minipage}[t]{0.48\textwidth}
\centering
\includegraphics[width=\linewidth]{issf_rifle.png}
\caption{ISSF 10\,m Air Rifle target \cite{issf2023_rules}}
\label{fig:issf_rifle_target}
\end{minipage}\hfill
\begin{minipage}[t]{0.48\textwidth}
\centering
\includegraphics[width=\linewidth]{issf_pistol.png}
\caption{ISSF 10\,m Air Pistol target \cite{issf2023_rules}}
\label{fig:issf_pistol_target}
\end{minipage}

\end{figure}

\begin{table}[ht]
\centering

\begin{minipage}[t]{0.48\textwidth}
\centering
\caption{ISSF 10\,m Air Rifle target ring diameters \cite{issf2023_rules}}
\label{tab:airrifle_rings}
\begin{tabular}{cc}
\hline
\textbf{Ring} & \textbf{Diameter (mm)} \\
\hline
10 & 0.5 \\
9  & 5.5 \\
8  & 10.5 \\
7  & 15.5 \\
6  & 20.5 \\
5  & 25.5 \\
4  & 30.5 \\
3  & 35.5 \\
2  & 40.5 \\
1  & 45.5 \\
\hline
\end{tabular}
\end{minipage}\hfill
\begin{minipage}[t]{0.48\textwidth}
\centering
\caption{ISSF 10\,m Air Pistol target ring diameters \cite{issf2023_rules}}
\label{tab:airpistol_rings}
\begin{tabular}{cc}
\hline
\textbf{Ring} & \textbf{Diameter (mm)} \\
\hline
10 & 11.5 \\
9  & 27.5 \\
8  & 43.5 \\
7  & 59.5 \\
6  & 75.5 \\
5  & 91.5 \\
4  & 107.5 \\
3  & 123.5 \\
2  & 139.5 \\
1  & 155.5 \\
\hline
\end{tabular}
\end{minipage}

\end{table}

Figure~\ref{fig:issf_rifle_target} illustrates the ISSF 10\,m air rifle paper target, while
Fig.~\ref{fig:issf_pistol_target} shows the corresponding 10\,m air pistol target. The
official ring diameters for both targets are listed in
Table~\ref{tab:airrifle_rings} and Table~\ref{tab:airpistol_rings}, respectively.

For the 10\,m air rifle paper target, the innermost ring (ring~10) is reduced to a small
white dot at the center of the target with a diameter of 0.5\,mm, as shown in
Fig.~\ref{fig:issf_rifle_target}. The ring diameters increase arithmetically by 5\,mm up to
the outermost ring (ring~1), which has a diameter of 45.5\,mm, as summarized in
Table~\ref{tab:airrifle_rings}. The inner black region extends from ring~10 to ring~4, with
a total diameter of 30.5\,mm.

The 10\,m air pistol paper target is larger, as shown in Fig.~\ref{fig:issf_pistol_target}.
The innermost ring (ring~10) has a diameter of 11.5\,mm, while the outermost ring (ring~1)
has a diameter of 155.5\,mm. The ring diameters increase arithmetically by 16\,mm, as listed
in Table~\ref{tab:airpistol_rings}. The inner black region extends from ring~10 to ring~7,
with a diameter of 59.5\,mm.

The dimensions summarized in Fig.~\ref{fig:issf_rifle_target} and
Table~\ref{tab:airrifle_rings} highlight the extremely small scale of the air rifle target,
which directly motivates the accuracy analysis presented in the following section.

The continuous radial scoring model can be written as
\begin{equation}
  s_c(r) = 11.0 - k\,r ,
  \label{eq:continuous_score}
\end{equation}
where $r$ is the radial distance between the shot center and the target center, and $k$ is a
scale factor and corresponds to the inverse radius increment
\begin{equation}
    k = \frac{2}{\Delta d_{1.0}}
\end{equation}
between adjacent rings.

ISSF decimal scoring applies rounding and clipping to the continuous scoring model with the formula
\begin{equation}
  s(r) =
  \mathrm{clip}\!\left(
    \frac{\left\lfloor 10\, s_c(r) \right\rfloor}{10},
    0,
    10.9
  \right),
  \label{eq:issf_decimal_score}
\end{equation}
where $\lfloor \cdot \rfloor$ denotes the floor operator and
$\mathrm{clip}(x,a,b)=\min(\max(x,a),b)$ limits the score to the valid ISSF range.

\section{Accuracy Requirements}

The scoring scale resolution and the smallest paper target size impose accuracy requirements for any scoring system.
Due to the small dimension of the 10m-air-rifle paper target, the required accuracy goes into the sub milimeter range.

To understand the accuracy requirements, we analyse the measurements relative to one interval of constant score. The length of the interval
is given by

\begin{equation}
  \Delta r_{0.1} = 0.05 \cdot \Delta d_{1.0}
\end{equation}

and corresponds to the increment in radial distance that produces a change of 0.1 in the continuous score.

The measurement error is modeled as a uniform random variable in the interval $[-E,\,E]$.
Let $R$ denote the ground-truth radial position within a decimal scoring interval and let
$\varepsilon$ denote the measurement error. The measured value is $M = R + \varepsilon$.
Assuming $R$ and $\varepsilon$ are independent, the probability density function (pdf) of $M$
is given by the convolution of the pdfs of $R$ and $\varepsilon$, as illustrated in
Fig.~\ref{fig:measurement_pdf}.

\begin{figure}[htbp]
  \centering
  \includegraphics[width=0.85\linewidth]{radial_error_pdf.png}
  \caption{Bounded-error model: uniform ground-truth position within a scoring interval and uniform measurement error, and the resulting measurement pdf obtained by convolution.}
  \label{fig:measurement_pdf}
\end{figure}

Regardless of how small the error is, there will always be a non-zero probability for incorrect scoring by 0.1.
The probability of correctly scoring a bullet is given by the integral of the pdf of the measurement within the scoring interval,
which is

\begin{equation}
  P = 1 - \frac{E}{2*\Delta r_{0.1}}
\end{equation}.

Solving for the accuracy E, we obtain:

\begin{equation}
  E = 2\,\Delta r_{0.1}\,(1 - P)
\end{equation}.

\begin{table}[htbp]
\centering
\caption{Required measurement accuracy $E$ as a function of the probability $P$ of correct decimal scoring for a 10\,m air rifle target ($\Delta r_{0.1} = 0.25$\,mm).}
\label{tab:accuracy_vs_probability}
\begin{tabular}{cc}
\hline
\textbf{Probability $P$ (\%)} & \textbf{Required accuracy $E$ (mm)} \\
\hline
80  & 0.10 \\
85  & 0.075 \\
90  & 0.050 \\
95  & 0.025 \\
99  & 0.005 \\
\hline
\end{tabular}
\end{table}


Table~\ref{tab:accuracy_vs_probability} shows that achieving high probabilities of correct
decimal scoring requires increasingly strict bounds on the radial measurement error.
For the 10\,m air rifle target, a probability of 90\% already requires a measurement
accuracy better than 0.05\,mm, while probabilities above 95\% demand accuracies below
0.03\,mm. These values illustrate the high accuracy requirements imposed by ISSF
decimal scoring and motivate the use of precise geometric localization methods in
automatic paper target scoring systems.

\section{Classical Computer Vision Approaches}
\label{sec:classical_methods}

Classical computer vision approaches to paper target scoring typically rely on explicit image
segmentation and geometric reasoning. In contrast to learning-based methods, they aim to
recover physically meaningful quantities---such as the target center, ring radii and bullet-hole
centers---and then compute the score by applying the official scoring rules. These pipelines can
achieve high localization accuracy under controlled imaging conditions, but they often become
fragile when lighting varies, the background is cluttered, or bullet holes overlap.

\subsection{Geometric scoring with perspective correction and contour-based circle fitting}
A representative geometric pipeline is proposed by Yu \textit{et al.}~\cite{yu2023_contourcircle},
who score paper targets using only classical image processing and circle fitting. The method first
computes a projective transformation to obtain a fronto-parallel view of the target, compensating
for camera viewpoint. After perspective correction, the outer contour of the target is extracted and
a circle is fitted to estimate the target center and scale. Bullet holes are then isolated via
thresholding and morphological operations, and connected components are extracted. For each
bullet-hole region, contours are computed and a circle model (e.g., minimum enclosing circle or
least-squares circle fitting) is fitted to estimate the bullet-hole center. Finally, the score is computed
from the distance between the bullet-hole center and the target center using the known ISSF ring
geometry.

The main strength of this approach is that it reconstructs the target geometry and estimates
centers via contour-based fitting, which is well-aligned with ISSF-compliant scoring. However, it
assumes that segmentation is reliable, which is typically achieved under controlled backgrounds and
stable lighting. As a result, performance degrades when bullet holes have bad contrast against the
paper, lighting changes, or when bullet holes overlap, since the connected components
no longer correspond to individual impacts.



\subsection{Segmentation of overlapping bullet holes using watershed and circle fitting}
Overlapping impacts are a major failure mode for connected-component-based
pipelines. Liu \textit{et al.}~\cite{liu2023_watershed} address this scenario using an improved
watershed segmentation strategy. Their pipeline applies grayscale preprocessing and morphological
operations, followed by area-based filtering to suppress non-hole regions such as ring markings and
text. Overlapping bullet holes are then separated using a distance-transform-driven watershed method,
and boundaries are extracted. For each separated region, least-squares
circle fitting is used to estimate individual hole centers and radii.

This work demonstrates that, when the preprocessing is successful, classical segmentation can separate
strongly overlapping holes and still recover geometrically meaningful centers. Nonetheless, the method
depends on threshold choices and assumptions about bullet-hole appearance (e.g., expected region
areas). In addition, the paper focuses on cropped regions around overlapping holes rather than full
target images, so it does not fully address end-to-end robustness in a complete target scoring system.



\subsection{Temporal differencing and Hough-based circle detection for live scoring}
A complementary classical strategy for live scoring is to exploit temporal information. Li
\textit{et al.}~\cite{li2022_imagesubtraction_hough} propose a method that processes pairs of images
captured before and after each shot. After geometric correction (homography), the pre-shot image is
subtracted from the post-shot image, only keeping the pixels that changed.
The pipeline then applies binarization, morphological cleaning and edge detection, followed by a
generalized Hough transform to detect the circular bullet-hole contour and estimate its center. A Hough
transform can also be used to localize a reference ring or the target center, enabling score computation
from the measured radial distance.

The main advantage of temporal differencing is facilitating the segmentation task. 
However, it requires a stable camera and assumes that
exactly one new shot occurs between consecutive frames; camera motion or multiple new impacts can
break the subtraction assumption. Despite these constraints, temporal differencing is attractive as an
optional component in real-time systems where images are captured continuously and the shot timing is
known.

\subsection{Discussion}
Overall, classical methods are attractive because they compute scores from explicit target geometry
and can be computationally lightweight. Their limitations arise primarily from the sensitivity of
segmentation to illumination changes, background clutter and target mounting frames (like a metal frame).
In addition, classical pipelines are typically parameter-heavy, relying on multiple thresholds
and heuristic parameters (e.g., intensity thresholds, morphological kernel sizes and area constraints)
that must be carefully tuned for each imaging setup. These observations motivate hybrid designs
that preserve geometric scoring and center estimation, while improving robustness through stronger
detection mechanisms (e.g., learning-based detectors or temporal differencing when available).

\section{Deep Learning-Based Approaches}
\label{sec:deep_learning_methods}

Recent work has explored the use of deep learning, particularly convolutional neural networks
(CNNs), for automatic detection of bullet holes on paper targets. These approaches typically frame
the problem as an object detection or instance segmentation task, with the goal of localizing bullet
holes under challenging imaging conditions.

A representative study is presented by Butt \textit{et al.}~\cite{butt2024_bulletholes_yolov8}, who apply the
YOLOv8 object detection framework to detect bullet holes on indoor shooting targets. The authors
report high detection performance, with mean Average Precision (mAP) values exceeding 95\% under
controlled indoor conditions. A key advantage of the YOLOv8-based approach is its computational
efficiency: inference times are on the order of a few milliseconds per
image, making the method suitable for real-time or near-real-time applications. The study
demonstrates that CNN-based detectors can robustly identify bullet holes despite moderate lighting
variations and background clutter, outperforming classical threshold-based methods in detection
reliability.

However, the detected bullet holes are represented by bounding boxes, and the paper does not
address sub-region analysis or precise center estimation beyond the detector output. Consequently,
while detection accuracy is high, the geometric precision required for ISSF-compliant decimal
scoring is not directly achieved.

A related line of work is presented by Ahmed \textit{et al.}~\cite{ahmed2023_maskrcnn_resnet}, who employ a
Mask R-CNN architecture with a ResNet-50 backbone to detect and segment bullet holes on outdoor
shooting targets. This work demonstrates the robustness of deep learning methods in more adverse
conditions, including variable illumination, weather effects and non-uniform backgrounds. The use
of instance segmentation allows the network to recover the approximate shape of bullet holes rather
than only bounding boxes, which is advantageous for visual interpretation and robustness.

Despite this increased robustness, the reported localization accuracy remains limited by the
resolution of the segmentation masks and the network output grid. The study focuses primarily on
detection and segmentation performance metrics and does not explicitly analyze the accuracy of the
estimated bullet-hole centers relative to the precision required for ISSF decimal scoring. In
addition, overlapping bullet holes are not explicitly addressed, and the segmentation results are
evaluated mainly on isolated impacts.

Overall, deep learning approaches offer significant advantages in terms of robustness to illumination
changes, background clutter and target wear, and they enable fast, real-time detection on low-cost
hardware. However, accurately estimating bullet-hole center coordinates with sub-millimeter
precision remains a major challenge. The detector outputs (bounding boxes or coarse masks) are not
sufficiently precise for high-accuracy scoring on their own, and none of the reviewed deep learning
studies explicitly address the problem of overlapping bullet holes. These limitations motivate hybrid
approaches that combine CNN-based detection for robustness with classical geometric refinement for
precision.

\section{Discussion and Summary}
\label{sec:chapter2_summary}

This chapter reviewed the physical and regulatory constraints imposed by ISSF paper targets and
analyzed their implications for automatic scoring systems. Due to the small dimensions of the
10\,m air rifle target and the use of decimal scoring, even small localization errors can lead to
incorrect scores. A simple bounded-error model was introduced to illustrate how scoring accuracy
depends directly on the measurement error, highlighting the high accuracy requirements
imposed by ISSF rules.

The state of the art in computer-vision-based target scoring was then reviewed. Classical image
processing approaches rely on geometric modeling and can achieve high localization
accuracy under controlled conditions. However, these methods are sensitive to illumination
changes, background clutter and overlapping bullet holes, and they often require
careful tuning of multiple heuristic parameters. In contrast, deep learning-based approaches
demonstrate strong robustness and real-time performance, particularly for bullet-hole detection,
but their outputs are typically limited to bounding boxes or coarse segmentation masks (Mask R-CNN), which do
not provide the geometric precision required for ISSF-compliant decimal scoring.

Overall, the literature shows that no single approach fully satisfies both the robustness and
precision requirements of automatic ISSF paper target scoring. These observations motivate a
hybrid system design that combines the robustness of learning-based detection with the precision
of classical geometric refinement. The following chapter presents the proposed system architecture
based on this hybrid approach.
